\section{Complementary Learning Systems}

One influential perspective from cognitive science and machine learning is the complementary learning system view, which separates fast adaptation from slower long-term consolidation~\cite{kumaran2016learning}. In its classical form, this view argues that a learning system benefits from maintaining more than one memory mechanism: a fast component that can react quickly to new information, and a slower component that integrates information gradually into a stable long-term store.

This perspective is relevant to lifelong model editing because the editing problem naturally contains both requirements. Newly arriving facts often need to affect system behavior quickly, especially in environments where knowledge changes cannot wait for heavy re-training. At the same time, a model that only stores recent updates in a transient way is not sufficient, because accepted edits should eventually be consolidated into a more stable representation. Complementary learning therefore offers a useful conceptual background for separating immediate responsiveness from durable integration, and for explaining why purely online insertion and durable long-term storage may place different demands on the system during deployment.

In the thesis, this idea is used as background rather than proof of a finished mechanism. It does not by itself determine how a model should be edited, where new facts should be stored, or how consolidation should be triggered. What it provides is a principled explanation for why one should not assume that a single mechanism must handle both rapid update and long-term stability equally well. This is especially relevant in the present research, where the strongest direct rewriting methods and the most operationally convenient methods do not coincide empirically, and where that mismatch motivates the search for a complementary rather than purely uniform editing architecture.
